Clausen and Scholze \cite{Scholze} introduced \emph{light condensed sets}. One
guiding intuition is that a compact Hausdorff space should be resolved by
profinite ones: it admits a surjection from a profinite space, and indeed a
semi-simplicial resolution by such spaces \cite{dyckhoff76-1}, along which its
cohomology can be computed.
This intuition is captured by using light condensed
sets, which are defined as the sheaves on the site of light profinite sets, i.e.\ of sequential limits of finite sets, for
the topology generated by finite jointly surjective families. Light condensed
sets form a topos, and light condensed anima form a higher topos; both carry an
internal language. \emph{Synthetic Stone duality} (SSD), introduced in
\cite{synthetic-stone-duality}, is an attempt to axiomatise what can be stated
and proved internally about light condensed anima: it extends homotopy type
theory by four axioms abstracted from properties of this higher
topos.\footnote{These axioms are consistent: a model has been constructed in
recent work of Coquand, H\"ofer and Sattler \cite{coquand-hofer-sattler},
building on \cite{CRS21}, and this is done in a constructive metatheory.} In SSD one reasons from these axioms alone and the
site never appears. Stone spaces are the types that are merely sequential limits
of finite types, overtly discrete types the ones that are merely sequential
colimits of finite types, and compact Hausdorff spaces the quotients of Stone
spaces by closed equivalence relations --- the intuition above, turned into a
definition. These classes are closed under $\Sigma$-types, so that a family of
Stone spaces indexed by a Stone space has a Stone total space; this is what lets
the whole development below be carried out for families rather than for single
spaces. Moreover no continuity condition is ever imposed or verified, since every
map is automatically continuous \cite{synthetic-stone-duality}.

In this paper we give an internal derivation of results of Dyckhoff
\cite{dyckhoff76,dyckhoff76-1} on the cohomology of compact Hausdorff spaces,
which in the condensed setting are proved externally \cite[Theorem
3.2]{Condensed}. The central point is that $H^n(S,\Z) = 0$ for all $n > 0$ when
$S$ is a Stone space, the internal counterpart of a profinite set;
everything else follows from this by a \v{C}ech argument. A key idea in the proof
of this vanishing is due to David W\"arn, who used it in the setting of synthetic
algebraic geometry to show that $H^n(X,\mathcal{O}_X) = 0$ for $n > 0$ and $X$
affine (unpublished). The two settings are closely analogous: Stone spaces
correspond to affine schemes and compact Hausdorff spaces to separated schemes,
the latter being singled out here by the requirement that their identity types be
closed, as separated schemes are by a closedness condition on the diagonal. In \cite[Section 5]{synthetic-stone-duality} this was carried
out in degree $1$ and with constant coefficients $\Z$: it was shown that
$H^1(S,\Z) = 0$ for $S$ Stone, that $H^1(X,\Z)$ for $X$ compact Hausdorff can be
computed by \v{C}ech cohomology, and that $H^1(\I,\Z) = 0$ for the unit interval
$\I$. Here we prove the general form of these results, with coefficients that may
vary over the space. Writing $\AbGroup_{\ODisc}$ for the type of overtly discrete
abelian groups --- equivalently, by \Cref{overtly-discrete-countably-presented},
of countably presented abelian groups --- we obtain:

\begin{itemize}
\item for $S$ Stone and $A : S \to \AbGroup_{\ODisc}$, we have $H^n(S,A) = 0$ for
  all $n \geq 1$ (\Cref{vanishing-cohomology-stone}), and likewise $H^1(S,G) = 0$
  for a dependent overtly discrete group $G$ which need not be abelian
  (\Cref{vanishing-first-stone});
\item for $X$ compact Hausdorff and $q : S \twoheadrightarrow X$ a surjection with
  $S$ Stone, writing $S(x)$ for the fiber $q^{-1}(x)$, the group $H^n(X,A)$ is
  naturally isomorphic to the $n$-th cohomology group $\check{H}^n(X,S,A)$ of the complex
  \[\Pi_{x:X}A(x)^{S(x)} \to \Pi_{x:X}A(x)^{S(x)^2} \to
    \Pi_{x:X}A(x)^{S(x)^3} \to \cdots\]
  (\Cref{cech-and-sheaf-agree});
\item the unit interval is acyclic (\Cref{cohomology-I}): every map from $\I$ to
  an overtly discrete type is constant, $H^1(\I,G) = 0$ for any overtly discrete group $G$ and $H^n(\I,A) = 0$ for
  any $n \geq 1$ and overtly discrete abelian group $A$.
\end{itemize}

The theorems of Dyckhoff at issue are the following. In \cite{dyckhoff76-1} he
constructs, for any space $X$, a semi-simplicial resolution $P_\bullet X \to X$ by
extremally disconnected spaces,\footnote{These are the projective objects among
compact Hausdorff spaces, by a theorem of Gleason \cite{gleason58}. In our
setting the resolution is instead by arbitrary Stone spaces, and no projectivity
is needed: the higher cohomology of any Stone space already vanishes
(\Cref{vanishing-cohomology-stone}).} and shows that
for $X$ paracompact Hausdorff the cohomology $H^\ast(X,S)$ with coefficients in an
abelian sheaf $S$ is that of the complex $\mathrm{Hom}_X(P_\bullet X, S)$. It is
\cite[Theorem 3.11]{dyckhoff76} that Scholze appeals to in \cite[Theorem
3.2]{Condensed} in order to identify, for $X$ compact Hausdorff, the cohomology of
$X$ computed in the topos of condensed sets with the cohomology computed from a simplicial hypercover by
profinite sets. Our results are
the internalisation of this step, since $H^n(X,A)$ is by definition
cohomology computed in the topos. The statements themselves are classical; what
is new here is that they are obtained internally, by elementary reasoning about
towers of finite types in a language where a compact Hausdorff space is simply a
type, rather than through projective resolutions and extremally disconnected
covers. They are also more general than the ones we started from, since the
coefficients may be a family $A : X \to \AbGroup_{\ODisc}$ varying over the
space --- which is the form in which the internal language states such things in
any case, as internal quantification ranges over families.

That this can be carried out at all depends on the higher types of homotopy type
theory. Cohomology has a direct definition there, $H^n(X,A) =
\propTrunc{\Pi_{x:X}\K(A(x),n)}_0$, with no resolutions and no derived functors,
and it can be used to prove theorems of ordinary algebraic topology
\cite{hott,evan-master-thesis}. The higher types also let us define a \emph{shape}
modality, by nullifying the interval, that is, by localising at the map $\I \to
\unit$ \cite{modalities}; this is how Shulman
\cite{shulman-Brouwer-fixed-point} obtains the shape modality of real-cohesive
homotopy type theory, as nullification at $\R$. In synthetic Stone duality it is
studied by Rashid \cite{rashid-borsuk-ulam}, who uses it to characterise $\R$ in
terms of the modal unit of the topological circle, and to prove the
two-dimensional Borsuk--Ulam theorem. What \Cref{cohomology-I} provides
is a large supply of modal types for it. Overtly discrete types are $\I$-local,
and so are $\B G$ for $G$ an overtly discrete group and $\K(A,n)$ for $A$ an
overtly discrete abelian group (\Cref{delooping-I-local}); more generally, every
Postnikov complete type whose set of connected components and homotopy groups are
overtly discrete is $\I$-local
(\Cref{postnikov-I-local}). This needs the vanishing of $H^n(\I,A)$ in all
degrees and for overtly discrete coefficients, which is what is proved here: in
\cite{synthetic-stone-duality} only the case $n = 1$ with coefficients $\Z$ was
available. It is likewise the vanishing of $H^1(\I,G)$ for $G$ overtly discrete
and not necessarily abelian that implies that $\B G$ is $\I$-local, which is what justifies
the effort spent on non-abelian cohomology in degree $1$.

The technical engine behind these computations is a form of Scott continuity,
which is suggested by the work of Barton and Commelin \cite{bc24}, where it is
one of the axioms, and which we develop in
\Cref{scott-continuity-for-cohomology}. In its general form
(\Cref{eine-scott-continuity-stone} and \Cref{scott-continuity-chaus}), which
was suggested to us by Reid Barton, it says
that the functor $\Pi$, sending a compact Hausdorff space $C$ together with a
family $I : C \to \ODisc$ to $\Pi_{x:C}I(x)$, commutes with sequential colimits.
Two special cases are used constantly: $\mathrm{colim}_i \Pi_{x:S}I_i(x) =
\Pi_{x:S}\,\mathrm{colim}_i I_i(x)$ for $S$ Stone
(\Cref{scott-continuity-right}), and $\mathrm{colim}_k(S_k \to I) = (\lim_k S_k
\to I)$ for $I$ overtly discrete (\Cref{scott-continuity-left}). The name is
justified by the following consequence: any map $(\Pi_{k:\N}S_k) \to I$ from a
countable product of Stone spaces into an overtly discrete type merely factors
through a finite sub-product $\Pi_{k<n}S_k$. Together with Tychonov's theorem and
its dual --- a product of overtly discrete types indexed by a Stone
space is overtly discrete (\Cref{tychonov-dual-stone}) --- and with the closure
properties of \Cref{stone-spaces} and \Cref{overtly-discrete-types}, this allows
us to verify \emph{all} their axioms in SSD (\Cref{barton-commelin}), of which
only some were obtained in \cite{synthetic-stone-duality}. Their theory is itself
an axiomatisation of the internal language, postulating at the level of sets two
classes of types with good closure properties which interact in this way; so what
is at stake here is not an internalisation but a comparison between two synthetic
accounts of the same situation. What it gives us is that everything proved in
their system is available in SSD, with their two classes interpreted as the
overtly discrete types and the compact Hausdorff spaces --- and, since SSD is formulated
in homotopy type theory rather than at the level of sets, that this can be
combined with the higher types on which the definition of cohomology rests.

Not everything can be obtained this way. Central to the work of Clausen and
Scholze is the internal projectivity of the free abelian group $\Z[\Noo]$ on the
generic convergent sequence, and W\"arn \cite{warn2024} has shown that this
property is not stable under base change to a slice. So ``$\Z[\Noo]$ is a
projective abelian group'' does not hold in the internal language, and the
property Clausen and Scholze use has to be established externally, in a
classical metatheory.

\paragraph{Organisation.} \Cref{scott-continuity-for-cohomology} is
preceded by generalities on Stone spaces (\Cref{stone-spaces}) and overtly
discrete types (\Cref{overtly-discrete-types}), and is followed by
\Cref{barton-commelin}, where Scott continuity is extended to compact Hausdorff
spaces so as to obtain the Barton--Commelin axioms; a reader interested only in
the cohomological results may skip that last subsection. The cohomology proper
begins with non-abelian coefficients over Stone and compact Hausdorff spaces
(\Cref{first-cohomology-stone-chaus}), continues with the vanishing of higher
cohomology for Stone spaces (\Cref{vanishing-stone}) and the comparison with
\v{C}ech cohomology for compact Hausdorff spaces (\Cref{cech-sheaf-agree-section}), and ends with the
acyclicity of the unit interval and its consequences for $\I$-local types
(\Cref{interval-acyclic}). The appendices contain results that hold in plain
homotopy type theory, without the axioms of synthetic Stone duality: the
comparison between non-abelian cohomology in degree $1$ and \v{C}ech cocycles
(\Cref{app-non-abelian-cech}), a criterion under which higher cohomology and
\v{C}ech cohomology agree (\Cref{app-abelian-cech}), and combinatorial results on
the cochain complex of an interval (\Cref{combinatorialIntervalSection}).

\paragraph{Prerequisites.} This paper is not self-contained: we assume
familiarity with \cite{synthetic-stone-duality}, whose axioms and basic results
we use freely. In \Cref{stone-spaces} and \Cref{overtly-discrete-types} we recall,
sometimes in slightly rephrased or extended form, only those results that are
needed later, and refer to \cite{synthetic-stone-duality} for their proofs.
